% Nouveau EGA V (cote 25): archive cote 25, site page 51 (PDF page 52).
% typescript. Draft compared: none.
% AI-transcribed by Claude Opus 5.5 (claude-opus-5-5, adaptive thinking, effort high), Message Batches API, batch ega25b (msgbatch_01A7n5YuXubUPNK1zsepi8RG), 2026-11-13; not yet checked by eye.
\scanpage{25/51}
\margin{\unsure{16}}
\margin{o--- c)} \inserted{Supposons cependant X lisse de dim 2, géom. intègre \unsure{propre} / k, et f une \uline{immersion}. Alors il résulte de 6.1. et 8.8. \inserted{et 8.14.} \unsure{que} si $Y^{\mathrm{sing}}_{\xi} \to \check{P}$ n'est pas gén. fini et inséparable, l'ens. des $\{\xi \in \check{P}$ tel que $Y_{\xi}$ soit \deleted{\unsure{géom.}} séparable sur $k(\xi)\}$ avec \uline{au plus 2} comp. irréd. géom. -- un compl. de codim $\geq 2$}
une quadrique de l'espace dual, relative à la forme duale ...).   Dans le
cas d'une surface non singulière \xout{dx} d'un espace projectif, cette situa-%
tion doit cependant être considérée comme exceptionnelle, cf nº suivant.

\margin{o---}
Nous allons examiner de plus près   maintenant le cas des surfaces
(le cas des courbes ne se pose évidemment pas du point de vue irréductibi-%
lité des sections hyperplanes \xout{\unsure{t}} ).

(NB  Je m'aperçois avec effroi que la quadrique n'a pas droit, dans
l'emploi que j'ai fait du mot ``\xout{f}réglé'', à être appelée réglée !  Cela doit
être en désaccord avec nos pères, et il faudra donc inventer un mot plus
adéquat pour la notion utilisée ici).

\uline{Proposition} 8.18.  Supposons \xout{X géom intègre} k alg clos, X \xout{intègre et}
\xout{et} intègre (resp. intègre et normale)\inserted{de dimension $d \geq 2$ et}
propre sur k, soit T une partie finie fermée de X telle que  X-T soit
lisse  et f $|$X-T  soit non ramifié.  Pour que l'ensemble des \hand{$\xi \in$} $\check{P}$  tels
que  $Y_{\hand{\xi}}$  soit \xout{\unsure{es}} géom irréd   (resp.  géom intègre)  de dimension  1
\unsure{fait}  un complémentaire de  codimension $\geq$ 2, il f et s que les deux
conditions suivantes soient satisfaites:

a)  Pour tout x\hand{$\in$}T, il existe une section hyperplane $Y_{\hand{\xi}}$   (\hand{$\xi \in$}P(k))
passant par x, de dimension \over{1}{$d-1$}, et qui soit irréductible.

b) \inserted{X'=} X-T est ``réglé'' (sic)  pour f, \uline{ou}  il existe  une section hyper
plane  $Y'_{\hand{\xi}}$  (\hand{$\xi \in$} $\check{P}$(k) )  de X',  \xout{singulière} qui soit de dimension \over{1}{d}\hand{-1},
singulière et  irréductible.

\deleted{Soit  en effet  Z  la  réunion  de  g($\overline{\underline{Y}'^{\hand{\mathrm{sing}}}}$), \xout{et} des hyperplans
de P  définis par les  f(x), x  T.  \xout{et}  Alors pour   P - Z,  Y   est
\uline{lisse} de dimension 1, \xout{\illegible{8}}
\xout{\illegible{10} au même \illegible{2} est lisse au dess}
\inserted{géom}
\xout{\unsure{des points de}}  et comme  il est connexe par 6.1. , il est géom irréd.}

\mspage{51}
